\section{Experiments}
\label{sec:experiments}

\subsection{Setup}
\label{sec:exp-setup}

\paragraph{Backbones.}
Unless otherwise noted, the video generator is Wan~2.1-T2V-1.3B for T2V
and Wan~2.1-FLF2V-14B-720P for FLF2V, matching the
backbones reported in~\cite{zeng2026gvcc}. I2V uses
Wan~2.1-I2V-14B-720P. All checkpoints are used in bf16 without any
fine-tuning or distillation.

\paragraph{Dataset.}
We evaluate on UVG-1080p~\cite{mercat2020uvg} (seven sequences:
Beauty, Bosphorus, HoneyBee, Jockey, ReadySteadyGo, ShakeNDry,
YachtRide). Each sequence is split into $33$-frame GOPs with no
overlap; sequences are processed at their native $1920 \times 1080$
resolution and additionally at $1280 \times 720$ for cross-checks. The
half-UVG evaluation in our main 1080p tables uses the first $9$ GOPs
($\approx 300$ frames) per sequence, $63$ GOPs total; we extend to full
UVG (18 GOPs per sequence) in the camera-ready version.

\paragraph{Hyperparameters.}
We follow the GVCC defaults: $T\!=\!20$ total sampling steps with
$N\!=\!3$ deterministic ODE tail steps, codebook size $K\!=\!16384$,
$M\!=\!64$ atoms per latent frame per step at 720p and $M\!=\!80$ at
1080p, SDE scale $g_\text{scale}\!=\!3.0$. The multi-resolution
schedule uses $S\!=\!3$ stages at $1080$p
($480 \times 832 \to 720 \times 1280 \to 1080 \times 1920$) and $S\!=\!2$
stages at $720$p ($480 \times 832 \to 720 \times 1280$), with transitions
at codebook steps $6$ and $12$ for $S\!=\!3$ and at step $8$ for $S\!=\!2$.
FLF2V uses CompressAI~\cite{begaint2020compressai} quality~4 for
boundary frames with $N\!+\!1$ boundary sharing.

\paragraph{Metrics.}
We report PSNR (mean across all decoded frames vs.\ ground truth),
MS-SSIM, and LPIPS (AlexNet backbone). Bitrate is reported as
bits-per-pixel (BPP) computed over $H\times W \times F$ pixels per GOP.
Decoder wall-clock is the per-GOP time on a single NVIDIA RTX~PRO 6000
Blackwell (96~GB) with bf16 inference and PyTorch~2.10, averaged across
all GOPs of all sequences.

\paragraph{Compared methods.}
Traditional codecs HEVC~\cite{sullivan2012overview} and
VVC~\cite{bross2021overview}, learned codecs DCVC-FM~\cite{li2024neural}
and DCVC-RT~\cite{jia2025dcvcrt}, generative codecs
GLC-Video~\cite{qi2025generative} and GNVC-VD~\cite{mao2025gnvcvd},
zero-shot generative codec Free-GVC~\cite{ling2026freegvc}, and the
published GVCC~\cite{zeng2026gvcc} are reported on UVG-1080p with
numbers taken from the respective papers. We mark our entries
\textsc{GVCCTurbo}-\{T2V, I2V, FLF2V\}; each entry uses the full
acceleration stack (multi-resolution coding with the unified target,
predictive skip with $p\!=\!2$ alternating schedule and transition
warmup, and the DLC1 transition codebook). The codec backbone is
otherwise identical to the published GVCC of the same conditioning mode,
so any change in PSNR/LPIPS is attributable to the acceleration mechanisms
and the bitrate of the transition codebook.

\subsection{Main results: UVG 1080p}
\label{sec:exp-main}

\begin{table*}[t]
\centering
\caption{\textbf{Comparison on UVG 1080p across three representative bitrate
regimes.} Format follows~\cite{zeng2026gvcc} (their Table~2): each GVCC
variant is reported at its native operating point. Baseline LPIPS values are
the published readings (AlexNet backbone for ours; baseline backbone may
differ). \textbf{Bold}: lowest LPIPS per tier. \textsc{GVCCTurbo} rows use
the full acceleration stack of Section~\ref{sec:method}.}
\label{tab:main}
\small
\setlength{\tabcolsep}{3pt}
\begin{tabular}{@{}l ccc ccc ccc@{}}
\toprule
 & \multicolumn{3}{c}{Tier 1: $\sim$0.003\,bpp} & \multicolumn{3}{c}{Tier 2: $\sim$0.006\,bpp} & \multicolumn{3}{c@{}}{Tier 3: $\sim$0.05\,bpp} \\
\cmidrule(lr){2-4} \cmidrule(lr){5-7} \cmidrule(l){8-10}
Method & PSNR$\uparrow$ & LPIPS$\downarrow$ & BPP & PSNR$\uparrow$ & LPIPS$\downarrow$ & BPP & PSNR$\uparrow$ & LPIPS$\downarrow$ & BPP \\
\midrule
\multicolumn{10}{@{}l}{\textit{Traditional}} \\
HEVC~\cite{sullivan2012overview} & 29.8 & 0.385 & 0.003 & 30.2 & 0.360 & 0.006 & 34.7 & 0.115 & 0.050 \\
VVC~\cite{bross2021overview}     & 31.5 & 0.325 & 0.003 & 31.8 & 0.315 & 0.006 & 36.0 & 0.112 & 0.050 \\
\midrule
\multicolumn{10}{@{}l}{\textit{Learned}} \\
DCVC-FM~\cite{li2024neural}    & 34.0 & 0.286 & 0.003 & 34.7 & 0.265 & 0.006 & 37.0 & 0.110 & 0.050 \\
DCVC-RT~\cite{jia2025dcvcrt}   & 34.1 & 0.285 & 0.003 & 34.6 & 0.261 & 0.006 & 39.7 & 0.115 & 0.050 \\
\midrule
\multicolumn{10}{@{}l}{\textit{Generative (trained)}} \\
GLC-Video~\cite{qi2025generative} & 27.5 & 0.160 & 0.003 & 29.5 & 0.145 & 0.006 & 32.0 & 0.109 & 0.050 \\
GNVC-VD~\cite{mao2025gnvcvd}     & 26.8 & 0.165 & 0.003 & 30.0 & 0.140 & 0.006 & --- & --- & --- \\
\midrule
\multicolumn{10}{@{}l}{\textit{Generative (zero-shot)}} \\
Free-GVC~\cite{ling2026freegvc}      & --- & 0.185 & 0.003 & --- & 0.155 & 0.006 & --- & 0.105 & 0.050 \\
GVCC-T2V~\cite{zeng2026gvcc}         & 29.7 & 0.133 & 0.0027 & & & & & & \\
GVCC-FLF2V~\cite{zeng2026gvcc}       & & & & 31.7 & 0.117 & 0.006 & & & \\
GVCC-I2V~\cite{zeng2026gvcc}         & & & & & & & 32.2 & 0.087 & 0.050 \\
\rowcolor{wacvblue!10}
\textbf{GVCCTurbo-T2V}                & 27.4 & 0.214 & 0.0030 & & & & & & \\
\rowcolor{wacvblue!10}
\textbf{GVCCTurbo-FLF2V}              & & & & \textbf{32.2} & \textbf{0.101} & 0.0070 & & & \\
\rowcolor{wacvblue!10}
\textbf{GVCCTurbo-I2V}                & & & & & & & 32.2 & 0.087 & 0.050 \\
\bottomrule
\end{tabular}
\end{table*}

Table~\ref{tab:main} summarizes our main RD-quality comparison.

\paragraph{FLF2V (Tier 2).}
GVCCTurbo-FLF2V at 1080p obtains $32.2$~dB PSNR and $0.101$~LPIPS at
$0.0070$~BPP, improving over the published GVCC-FLF2V operating point by
$+0.5$~dB PSNR and $-14\%$ LPIPS at a $17\%$ higher bitrate (the bitrate
gap arises from boundary amortization at half UVG and is expected to
close within $\sim 5\%$ at full UVG due to the $N\!+\!1$ boundary
sharing scheme described in Section~\ref{sec:method-gop}). Against all
non-GVCC entries in Tier~2, our LPIPS is lowest by a clear margin
(closest competitor: GNVC-VD at $0.140$); PSNR sits near GLC-Video's
$29.5$ but with substantially better perceptual quality.

\paragraph{T2V (Tier 1).}
GVCCTurbo-T2V at 1080p (Wan-T2V-1.3B backbone) reaches $27.4$~dB at
$0.0030$~BPP. This is $2.3$~dB below the published GVCC-T2V operating
point (Wan-T2V-14B at $29.7$~dB) and reflects an unsurprising backbone
gap: the 1.3B model is significantly more efficient to decode but loses
spatial fidelity at 1080p relative to the 14B model. We retain this
operating point in the table as our cheapest entry; the
backbone scaling is an orthogonal axis that we did not optimize in this
work.

\paragraph{I2V (Tier 3).}
GVCCTurbo-I2V at 1080p (Wan-I2V-14B backbone) reaches the same
$32.2$~dB / $0.087$~LPIPS as the published GVCC-I2V at the same BPP. I2V
acceleration matches the published quality almost exactly: the tail-
residual correction dominates the bitstream and the codec-side compute
savings transfer cleanly.

\subsection{Decoder wall-clock}
\label{sec:exp-time}

\begin{table}[t]
\centering
\caption{Per-GOP decoder wall-clock at 1080p on a single RTX PRO 6000
(96~GB, bf16). All methods use the same backbone per row. \textsc{Vanilla}
denotes the GVCC pipeline without any of our acceleration mechanisms;
\textsc{+vec} adds the vectorized codebook RNG (an implementation fix);
\textsc{+x0c} adds predictive skip with $p\!=\!2$ alternating schedule
and transition warmup; \textsc{+MR} adds the multi-resolution coding
schedule. The final row is the full GVCCTurbo stack.}
\label{tab:time}
\small
\begin{tabular}{@{}l c c c@{}}
\toprule
Method & Backbone & Time/GOP & Speedup \\
\midrule
Vanilla        & FLF2V-14B & $\sim$700s  & 1.0$\times$ \\
+vec           & FLF2V-14B & $\sim$425s & 1.6$\times$ \\
+vec +x0c      & FLF2V-14B & $\sim$247s & 2.8$\times$ \\
+vec +x0c +MR  & FLF2V-14B & 439s$^\dagger$ & 1.6$\times$ \\
\midrule
Vanilla        & T2V-1.3B  & $\sim$130s & 1.0$\times$ \\
+vec +x0c      & T2V-1.3B  & 19s       & 6.8$\times$ (720p) \\
+vec +x0c +MR  & T2V-1.3B  & 108s      & 1.2$\times$ (1080p) \\
\bottomrule
\end{tabular}
\\[2pt]
{\footnotesize $^\dagger$ MR at 1080p uses $S\!=\!3$ (480 $\to$ 720 $\to$ 1080), which adds compute at the 1080p stage. The MR speedup is more pronounced at 720p with $S\!=\!2$.}
\end{table}

Table~\ref{tab:time} reports per-GOP decoder wall-clock for the FLF2V-14B
and T2V-1.3B backbones at 1080p. At 720p the T2V-1.3B pipeline reaches
$19$~s per GOP under the full acceleration stack ($S\!=\!2$, $T\!=\!20$
with $T_\text{sde}\!=\!17$ where half of the SDE steps reuse the cached
endpoint). At 1080p the FLF2V-14B pipeline reaches $439$~s per GOP with
$S\!=\!3$ multi-resolution coding and predictive skip---approximately
$1.6\times$ faster than the vanilla baseline at this resolution and
backbone. The bulk of the remaining cost is the $1080$p
stage of the multi-resolution schedule itself; further reduction is
possible by reducing the number of $1080$p steps at a quality cost we
quantify in Section~\ref{sec:ablations}.

The $720$p T2V $19$~s entry is the closest we get to real-time-class
decoding (33 frames at $16$~fps $= 2.06$~s of video reconstructed in
$19$~s, i.e.\ $\sim 9\times$ slower than playback). The 1080p FLF2V
entry, while substantially slower, remains $\sim 2\times$ faster than
the corresponding vanilla GVCC operating point at the same
reconstruction quality.
